\section{Positional Encoding}
\label{sec:posenc}

The model supports two positional encoding strategies, configured via the
\texttt{pos\_type} field:

\subsection{Learnable Positional Embeddings (Default)}

The default configuration uses a \textbf{learnable embedding table}
$\mathbf{P} \in \mathbb{R}^{L_\text{ctx} \times d_\text{model}}$, where each
position $i$ has a trainable vector $\mathbf{p}_i$.  The positional signal is
added to the token embedding:
\[
\mathbf{h}_i^{(0)} = \mathbf{E}[x_i] + \mathbf{P}[i]
\]

This approach is simple and effective for fixed-length contexts.  The key
challenge arises during context extension (e.g., 256$\to$384): the new model
has more positions than the trained embedding table.

\textbf{Context extension via interpolation.}
When expanding context length, the positional embedding matrix is extended via
1D linear interpolation along the sequence dimension.  The old embeddings
$(L_\text{old}, d)$ are interpolated to $(L_\text{new}, d)$ using
\texttt{F.interpolate} with \texttt{align\_corners=True}.  Existing positions
retain near-identical values, while new positions receive smooth estimates.
This is applied during Stages~2, 4, and 6.

\subsection{Rotary Position Embeddings (RoPE)}

The model also implements RoPE as an alternative.  RoPE encodes relative
positions directly into the query and key vectors via rotation matrices:
\[
q'_i = \mathbf{R}_{\theta,i} \cdot q_i, \quad k'_j = \mathbf{R}_{\theta,j} \cdot k_j
\]
where $\mathbf{R}_{\theta,i}$ applies rotations at frequencies determined by
$\theta = 10000^{-2k/d}$.  RoPE has the advantage that context extension is
automatic (no learned positions to resize), but the current training uses
learnable embeddings to maintain consistency with the TinyStories baseline.
